% ltex: language=it
\documentclass[12pt]{article}

% Shared preamble from project root
\usepackage{lecture}
\usepackage{acronyms}

\begin{document}
\section{Da MLP a CNN}\label{sec:da_mlp_a_cnn} % (fold)
Hubel \& Wiesel scoprono la presenza, nella corteccia visiva del gatto, di due tipologie di neuroni:
\begin{itemize}
    \item \textbf{simple cells}: agiscono come feature detector locali, fornendo \textbf{selettività};
    \item \textbf{complex cells}: fondono (\red{pooling}) gli output di \textbf{simple cells} in un intorno, garantendo \textbf{invarianza};
\end{itemize}
\begin{note}[Neocognitron]
    È una delle prime reti neurali che cerca di modellare questo comportamento
\end{note}
Le \acf{cnn} introdotte da LeCun a partire dal $1998$.
Le principali differenze rispetto a \textbf{MLP} sono:
\begin{itemize}
    \item \textbf{processing locale}: i neuroni sono connessi solo \textbf{localmente} ai neuroni del livello precedente.
        Ogni neurone esegue quindi un'elaborazione locale, portando a una forte \textbf{riduzione del numero di connessioni}.
        \item \textbf{pesi condivisi}: i pesi sono \textbf{condivisi} a gruppi.
            Neuroni diversi dello stesso livello eseguono lo stesso tipo di elaborazione su porzioni diverse dell'input, portando a una forte \textbf{riduzione del numero di pesi}.
\end{itemize}
\begin{example}{}{}
    \begin{cimage}
        \includegraphics[width=.2\textwidth]{assets/20261123_133201.png}
    \end{cimage}
    Ciascuno dei $4$ neuroni a destra è connesso solo a $3$ neuroni del livello precedente, con pesi condivisi.
    In totale si hanno $12$ connessioni e $3$ pesi contro le $24$ connessioni e i $24$ pesi di una equivalente porzione di \textbf{MLP}.
\end{example}
\begin{tip}[Alternanza dei livelli di feature extraction e pooling]
\end{tip}
% section Da MLP a CNN (end)
\end{document}

